
\paragraph{Approximating the second moment.} As discussed in~\Cref{sec:sample-loss-lmo} in the main text ,
$\overline\Sigma_t$ cannot be directly estimated or even stored. 
This is because it depends on the  per-sample
gradients that backpropagation never materializes, and it is too large
to store. Two approximations remove the two obstacles.
\begin{itemize}[leftmargin=1.5em,itemsep=2pt,topsep=2pt]
  \item \emph{Batch collapse.} We replace the within-batch average
  $\frac1{n_s}\sum_ig_{s,i}g_{s,i}^\top$ by the outer product
  $g_sg_s^\top$ of the batch gradient $g_s=\vec(G_s)$. This substitution
  underlies the second moments of Adam~\cite{adam} and
  Shampoo~\cite{shampoo}, and it is exact at batch size one. Under it
  the weights $w_{t,s,i}$ sum within each batch to
  $w_{t,s}=(1-\beta_2)\beta_2^{t-s}$, so the collapsed estimate is a
  weighted average over batch gradients,
  \[
    \overline\Sigma_t\approx\E_t[gg^\top],
    \qquad\text{where}\qquad
    \E_t[\phi(G)]=\sum_{s\le t}w_{t,s}\,\phi(G_s).
  \]
  \item \emph{Kronecker factorization.} We assume the
  Kronecker-factored form of~\eqref{eq:pq-factor}, with
  $P,Q$ diagonal as in \Cref{alg:ours}. Since $Q\otimes P$ is unchanged
  by the rescaling $(P,Q)\mapsto(aP,a^{-1}Q)$, we fix
  $\norm{\diag P}_\infty=\norm{\diag Q}_\infty=1$ and let an additional scalar
  $\kappa_t>0$ carry the magnitude:
  \begin{equation}\label{e-diag-kron-second-mom-model}
    \E_t[gg^\top]\approx\kappa_t\,(Q\otimes P).
  \end{equation}
\end{itemize}

\subsection{Estimating the Weighted Moments}
\label{sec:weighted-estimation}
In LoRA training, backpropagation gives the factor gradients $G_A=B^\top G$
and $G_B=GA^\top$~\eqref{e-grad-wrt-factors} at no extra cost. Forming
the full gradient $G$, by contrast, would take a $d_{\mathrm{out}}\times
d_{\mathrm{in}}$ matmul and full-matrix storage per layer, exactly what
the low-rank structure avoids. We therefore estimate $P$ and $Q$ from the
factor gradients. \rob{$p$ and $q$ need to defined for this section, or at least a reminder}

\paragraph{Moments of the factor gradients.} Since
$\vec(G_A)=(I\otimes B^\top)\vec(G)$ and conjugation acts factorwise on
a Kronecker product, the second-moment
model~\eqref{e-diag-kron-second-mom-model} gives
\begin{equation}\label{eq:observable-moments}
  \Sigma_A:=\E_t\big[\vec(G_A)\vec(G_A)^\top\big] \approx\kappa_t\,Q\otimes C_A,
  \qquad
  \Sigma_B:=\E_t\big[\vec(G_B^\top)\vec(G_B^\top)^\top\big] \approx \kappa_t\,P\otimes C_B,
\end{equation}
\rob{These identities are only true if $A$ and $B$ do not change over iterations. We need to at least spell this out.}
where $C_A=B^\top PB$ and $C_B=AQA^\top$ are both the $r\times r$ matrices, and the second approximation in~\eqref{eq:observable-moments} follows by
transposing $G$.

\paragraph{Fitting the factors.} We recover the factors by fitting the
second-moment model to the observed moments~\eqref{eq:observable-moments}. Write $\tilde q=\kappa_t\,q$ and
$\tilde p=\kappa_t\,p$ for the scaled factors, 
\rob{It should be $q=\sqrt{\kappa_t}\,q$ and $\tilde p=\sqrt{\kappa_t}\,p$ ?}
absorbing the
magnitude that the normalization later removes, and let
$\tilde Q=\diag(\tilde q)$ and $\tilde P=\diag(\tilde p)$.
For a divergence $D$ between positive definite matrices, the fits are
\begin{equation}\label{eq:factor-fit}
  \widehat q
  =\argmin_{\tilde q>0}\;
  D\big(\Sigma_A,\;\tilde Q\otimes C_A\big),
  \qquad
  \widehat p
  =\argmin_{\tilde p>0}\;
  D\big(\Sigma_B,\;\tilde P\otimes C_B\big).
\end{equation}
The divergence resolves an overdetermined fit: the candidate family
$\{\tilde Q\otimes C_A\}$ has $d_{\mathrm{in}}$ parameters while
$\Sigma_A$ has on the order of $(rd_{\mathrm{in}})^2$ entries, and
similarly for the fit for $p$. Away from the exact model, no candidate
reproduces the observed moment, and $D$ decides how the mismatch is
weighed. We require two properties of $D$:
\begin{enumerate}[label=(\roman*),ref=(\roman*),leftmargin=2.2em,itemsep=1pt,topsep=2pt]
  \item\label{it:block-additive} \emph{additivity across blocks:}
  $D(\Sigma,S)=\sum_jD(\Sigma_j,S_j)$ whenever $\Sigma$ and $S$ are
  block diagonal with blocks $\Sigma_j$ and $S_j$ of matching sizes;
  \item\label{it:scalar-fit} \emph{a homogeneous scalar fit:} the best scalar multiple
  $c^\star(\Sigma,S)=\argmin_{c>0}D(\Sigma,cS)$ exists, is unique, and
  satisfies $c^\star(a\Sigma,S)=a\,c^\star(\Sigma,S)$ for $a>0$.
\end{enumerate}
Each fit takes the other factor as given ($C_A$ contains $P$, and
$C_B$ contains $Q$), so the fits are coupled. In practice, each is
posed with the stored estimate of the other factor. The next lemma
shows the coupling is sound.

\begin{lemma}[Consistency]\label{lem:consistency}
Let $D$ satisfy \ref{it:block-additive} and \ref{it:scalar-fit}, and let the
model~\eqref{e-diag-kron-second-mom-model} hold with equality. For any
diagonal $\widehat P\succ0$ and $\widehat Q\succ0$, the
fits~\eqref{eq:factor-fit} posed with $\widehat C_A=B^\top\widehat PB$
and $\widehat C_B=A\widehat QA^\top$ return
\[
  \widehat q=\kappa_t\,c^\star\big(C_A,\widehat C_A\big)\,\diag(Q),
  \qquad
  \widehat p=\kappa_t\,c^\star\big(C_B,\widehat C_B\big)\,\diag(P),
\]
so $\widehat q/\norm{\widehat q}_\infty=\diag(Q)$ and
$\widehat p/\norm{\widehat p}_\infty=\diag(P)$.
\end{lemma}
\begin{proof}
Under the model, \eqref{eq:observable-moments} makes $\Sigma_A$ block
diagonal with $j$th block $\kappa_t\,q_j\,C_A$, and the candidate
$\tilde Q\otimes\widehat C_A$ is block diagonal with $j$th block
$\tilde q_j\,\widehat C_A$. By \ref{it:block-additive} the objective
in~\eqref{eq:factor-fit} is
$\sum_jD(\kappa_t q_jC_A,\,\tilde q_j\widehat C_A)$, one term per
coordinate, and by \ref{it:scalar-fit} the $j$th term is minimized at
$\tilde q_j=\kappa_t\,q_j\,c^\star(C_A,\widehat C_A)$. The factor
$\kappa_t\,c^\star(C_A,\widehat C_A)$ is the same for every $j$, and
$\norm{\diag Q}_\infty=1$, so the normalization removes it. The claim
for $\widehat p$ follows with the roles of the two factors exchanged.
\end{proof}

\paragraph{Choosing the divergence.} The choice of divergence
determines the estimator. We consider two divergences that satisfy
\ref{it:block-additive} and \ref{it:scalar-fit}. \methodname{} uses the first.
\begin{itemize}[leftmargin=1.5em,itemsep=2pt,topsep=2pt]
  \item \emph{Log-determinant} (following KL-Shampoo~\cite{klshampoo}):
  $D(\Sigma,S)=\Tr(\Sigma S^{-1})-\log\det(\Sigma S^{-1})-d$, where $d$
  is the common dimension. It satisfies \ref{it:block-additive} since traces and
  log-determinants add across blocks, and \ref{it:scalar-fit} with scalar fit
  $c^\star(\Sigma,S)=\Tr(\Sigma S^{-1})/d$. Substituting
  $S=\tilde Q\otimes C_A$ into~\eqref{eq:factor-fit} separates the
  objective across coordinates, and the minimizers are
  \[
    T_Q(P)=\frac1r\diag\!\left(\E_t[G_A^\top C_A^{-1}G_A]\right),
    \qquad
    T_P(Q)=\frac1r\diag\!\left(\E_t[G_B C_B^{-1}G_B^\top]\right).
  \]
  The moments $\Sigma_A,\Sigma_B$ enter only through these products, so
  neither is ever formed, and the fits consume only the factor
  gradients and the $r\times r$ matrices $C_A$, $C_B$ that the
  direction step already uses.
  \item \emph{Von Neumann}:
  $D(\Sigma,S)=\Tr(\Sigma\log\Sigma-\Sigma\log S-\Sigma+S)$ also
  satisfies \ref{it:block-additive} and \ref{it:scalar-fit}, with scalar fit
  $c^\star(\Sigma,S)=\Tr(\Sigma)/\Tr(S)$. The same substitution yields
  the raw second moments of the factor gradients,
  \[
    \tilde q_j=\E_t\big[(G_A^\top G_A)_{jj}\big]/\Tr(C_A),
    \qquad
    \tilde p_i=\E_t\big[(G_B G_B^\top)_{ii}\big]/\Tr(C_B),
  \]
  with $C_A$, $C_B$ entering only through scalars that the
  normalization removes. This mirrors the correspondence established in
  KL-Shampoo between the von Neumann divergence and Adafactor-style
  moment matching~\cite{klshampoo}.
\end{itemize}
Both choices recover the factors under the model
(\Cref{lem:consistency}) up to a scaling factor. Where they differ is that the
log-determinant fit weighs the average by $C_A^{-1}$, where-as the von
Neumann does not.
